\section{Related Work}
Recent studies have benchmarked pathology foundation models across downstream predictive performance, transferability, and clinical tasks \citep{lei2026validation, boutaj2026from, vanderbilt2026goldmark}. With the growing number of pathology FMs, multi-FM fusion has emerged as a complementary direction, ranging from fixed aggregation such as concatenation and averaging to learnable gated and attention-based fusion. Other approaches explore adaptive ensemble paradigms, including logit-level fusion \citep{huang2026plugandplay}, group-based encoder selection \citep{quan2025gas}, and ensemble learning for precision oncology \citep{luo2025elf}. Recent methods further investigate unified or adaptive integration of multiple FMs, including Meta-Encoder, AdaFusion, Shazam, and FuseCPath \citep{gao2026metaencoder, xiao2025adafusion, lei2026shazam, yang2025fusecpath}. These methods primarily focus on learning fusion or selection mechanisms, but do not systematically analyze how representation composition affects individual FM task-level contributions.

Mixture-of-experts (MoE) approaches provide another paradigm for adaptive expert selection. Recent studies have explored multimodal MoE for immune biomarker prediction and structured routing based on region-graph optimal transport for WSI classification \citep{liu2026predicting,tian2026regiongraph}. These works highlight the potential of expert specialization and adaptive routing, but primarily study expert selection or structured routing rather than how the composition of foundation-model representations affects downstream fusion behavior.

Representation similarity has been widely studied for characterizing relationships among learned feature spaces, with centered kernel alignment (CKA) being a common measure \citep{cka}. Recent work has also analyzed pathology FM representations, reporting slide-dependent representation structure and differences between vision-only and vision-language models \citep{mishra2025comparing}. However, these studies mainly characterize individual representations rather than their task-conditioned interactions within multi-FM fusion. We therefore conduct a multi-perspective analysis of representation geometry, feature association, and task-level contribution, and examine how representation composition relates to fusion behavior.